\documentclass[11pt,a4paper]{article}

\usepackage[utf8]{inputenc}
\usepackage[T1]{fontenc}
\usepackage{lmodern}
\usepackage{geometry}
\usepackage{enumitem}
\usepackage{hyperref}
\usepackage[dvipsnames]{xcolor}
\usepackage{array}
\usepackage{soul}
\usepackage{amsmath}
\usepackage{amssymb}

\geometry{margin=2.2cm}
\hypersetup{
  colorlinks=true,
  linkcolor=black,
  urlcolor=blue,
  citecolor=black
}
\setlist[itemize]{leftmargin=1.2cm,itemsep=2pt,topsep=4pt}
\setlist[enumerate]{leftmargin=1.0cm,itemsep=4pt,topsep=6pt}
\setlength{\parindent}{0pt}
\setlength{\parskip}{6pt}
\renewcommand{\arraystretch}{1.2}

\newcommand{\mvp}[1]{{\setulcolor{ForestGreen}\ul{#1}}}
\newcommand{\str}[1]{{\setulcolor{BurntOrange}\ul{#1}}}

% Ukloni automatski naslov ("References") iz thebibliography okruženja
\makeatletter
\renewenvironment{thebibliography}[1]
     {\list{\@biblabel{\@arabic\c@enumiv}}%
           {\settowidth\labelwidth{\@biblabel{#1}}%
            \leftmargin\labelwidth
            \advance\leftmargin\labelsep
            \@openbib@code
            \usecounter{enumiv}%
            \let\p@enumiv\@empty
            \renewcommand\theenumiv{\@arabic\c@enumiv}}%
      \sloppy
      \clubpenalty4000
      \@clubpenalty \clubpenalty
      \widowpenalty4000%
      \sfcode`\.\@m}
     {\def\@noitemerr
       {\@latex@warning{Empty `thebibliography' environment}}%
      \endlist}
\makeatother

\begin{document}

\begin{center}
{\Large Univerzitet u Nišu\\Elektronski fakultet}\\[1.2cm]
{\LARGE \textbf{Predlog projekta iz predmeta Računarski vid}}\\[0.6cm]
{\Large \textbf{Uporedna analiza reprezentacija slike primenjenih na klasifikaciju saobraćajnih znakova pri kontrolisanim degradacijama ulaza}}\\[1.2cm]
{\large Student: Mihajlo Madić, broj indeksa 2119}\\[0.3cm]
{\large Niš, 2026.}
\end{center}

\vspace{0.6cm}

\section*{1. Naziv i kratak opis projekta}

\textbf{Naziv projekta:} \textit{Uporedna analiza reprezentacija slike primenjenih na klasifikaciju saobraćajnih znakova pri kontrolisanim degradacijama ulaza.}

Cilj projekta je sistematsko poređenje četiri različite reprezentacije~\cite{repr-learning} slike saobraćajnog znaka, od kojih su tri dobijene "klasičnim" metodima, a jedna je naučena -- metodama dubokog učenja. Date reprezentacije se koriste za izradu rešenja klasifikacije saobraćajnih znakova, nad GTSRB~\cite{gtsrb} (\textit{German Traffic Sign Recognition Benchmark}) skupom podataka, koji će biti korišćen za razvoj i evaluaciju rešenja. Ulaz sistema je isečak slike (eng. \textit{crop}), dimenzija između $15\times15$ i $250\times250$ piksela, koji sadrži tačno jedan saobraćajni znak, a izlaz je jedna od 43 klase znakova.

Suština projekta nije u postizanju najviše moguće tačnosti klasifikacije, GTSRB je zasićen skup podataka na kome savremene konvolucione mreže prelaze $99\%$~\cite{ciresan}, već u \textbf{analizi pod kojim uslovima je koja reprezentacija primenljiva}. Četiri posmatrane reprezentacije se razlikuju po više strukturnih osobina istovremeno (globalna nasuprot lokalnoj, sa očuvanim nasuprot odbačenim prostornim rasporedom, ručno projektovana tj. "klasična" nasuprot naučenoj), pa se očekuje da \mvp{redosled uspešnosti metoda zavisi od vrste degradacije ulazne slike}. Upravo ta \textit{interakcija} između reprezentacije i tipa degradacije predstavlja glavni rezultat projekta\footnote{Kroz ceo dokument su \textcolor{ForestGreen}{zelenom bojom} podvučeni elementi koji čine minimalni skup rezultata -- ono što autor smatra da mora biti realizovano da bi se projekat smatrao uspešnim. \textcolor{BurntOrange}{Narandžastom bojom} podvučen je \emph{prošireni skup ciljeva}, koji se realizuju samo ako vreme i resursi dozvole.}. 

Ovakva postavka problema je relevantna zato što daje informaciju pri odabiru date reprezentacije u implementaciji realnog sistema. U realnim okolnostima, odabir određene reprezentacije mora biti donešen u skladu sa nedostacima sistema u okviru koga se dati klasifikator koristi, a u toj odluci je najrelevantnija otpornost klasifikatora na onu vrstu degradacije koja u datoj primeni dominira, npr. zamućenje usled kretanja, šum pri slabom osvetljenju ili nepreciznost izlaza detektora.

\section*{2. Projektni tim}

\begin{tabular}{|>{\raggedright\arraybackslash}p{7cm}|>{\raggedright\arraybackslash}p{4cm}|}
\hline
\textbf{Ime i prezime} & \textbf{Broj indeksa} \\
\hline
Mihajlo Madić & 2119 \\
\hline
\end{tabular} \\

\newpage

\section*{3. Kategorija projekta}

Projekat pripada \textbf{tehnološkoj kategoriji}, sa \textit{eksperimentalno-analitičkim} karakterom.

Tehnološki deo je u implementaciji i praktičnom poređenju metoda obrađenih na časovima -- PCA (metod sopstvenih vektora, isti princip kao Eigenfaces~\cite{eigenfaces}), HOG deskriptora~\cite{hog}, metoda vreća vizuelnih reči nad SIFT deskriptorima~\cite{sift, bovw}, i konvolucione neuronske mreže -- ali u okviru jedinstvenog, kontrolisanog eksperimentalnog protokola. Za razliku od primera sa časova, gde se svaka od ovih metoda demonstrira izolovano i na sopstvenom skupu podataka, ovde se sve metode primenjuju na isti skup, uz isti pretprocesing, isti klasifikator i istu metodologiju evaluacije, čime se razlike u rezultatima mogu pripisati isključivo samoj reprezentaciji.

Analitički deo je u tome što se ne izveštavaju samo skalarne vrednosti tačnosti, već i ponašanje metoda pod kontrolisanim degradacijama ulaza, uz unapred zapisana očekivanja i njihovu naknadnu proveru.

\section*{4. Detaljniji opis projekta i očekivani rezultati}

\subsection*{4.1 Kontekst -- mesto projekta u širem sistemu}

Kompletan sistem za prepoznavanje i razumevanje saobraćajne signalizacije iz video snimka, snimljenog kamerom uperenom u smeru kretanja vozila, u literaturi se standardno razlaže na fazu detekcije i fazu klasifikacije~\cite{tsr-survey, gtsdb}. Radi preciznijeg razgraničenja opsega ovog projekta, ta podela je ovde dodatno raščlanjena na pet modula:

\begin{center}
\begin{tabular}{|c|p{7.6cm}|p{4.6cm}|}
\hline
\textbf{Modul} & \textbf{Uloga} & \textbf{Status u ovom projektu} \\
\hline
M1 & \textbf{Detekcija / lokalizacija} -- pronalaženje kandidatnih regiona koji sadrže znak u punom kadru & \textit{Van opsega} \\
\hline
M2 & \textbf{Praćenje} -- propagacija detektovanog regiona kroz naredne kadrove pomoću optičkog toka, radi uštede računskih resursa & \textit{Van opsega} \\
\hline
M3 & \textbf{Prepoznavanje (klasifikacija)} -- dodeljivanje klase izdvojenom isečku & \textbf{\mvp{Predmet ovog projekta}} \\
\hline
M4 & \textbf{Vremenska agregacija} -- objedinjavanje odluka po jednom fizičkom znaku kroz više kadrova & \textit{Van opsega} \\
\hline
M5 & \textbf{Semantičko tumačenje} -- prevođenje prepoznatih znakova u ograničenja i preporuke za vozilo & \textit{Van opsega} \\
\hline
\end{tabular}
\end{center}

\textbf{Ovaj projekat realizuje isključivo modul M3, i to nad isečcima čije su granice date anotacijom skupa podataka, a ne izlazom detektora.} Moduli M1, M2, M4 i M5 se u ovom dokumentu opisuju samo radi kontekstualizacije -- oni nisu deo predmeta rada, ne implementiraju se, i ne evaluiraju se. Ovo ograničenje opsega je namerno i proizilazi iz opsega samog predmeta i vremenskih resursa; obrazloženje i pravac daljeg razvoja dati su u Dodatku A.

Jedan element projekta je ipak svesno oblikovan tako da služi budućem povezivanju sa modulom M1: \str{merenje osetljivosti klasifikatora na pomeraj i promenu razmere graničnog okvira, nad GTSDB skupom, koji sadrži kadrove punih scena}~\cite{gtsdb} (opisano detaljnije u odeljku 6.2). Realan detektor nikada ne daje isečke tako precizno centrirane i uokvirene kao ručna anotacija, pa ovo merenje kvantifikuje koliko bi tačnosti bilo izgubljeno kada bi se modul M3 spojio sa modulom M1.

\subsection*{4.2 Formulacija problema}

\begin{itemize}
  \item \textbf{Ulaz:} slika $I \in \mathbb{R}^{h \times w \times 3}$, isečak koji sadrži tačno jedan saobraćajni znak; $h, w \in [15, 250]$.
  \item \textbf{Izlaz:} oznaka klase $\hat{y} \in \{1, \dots, 43\}$.
  \item \textbf{Zadatak:} naučiti preslikavanje $I \mapsto \hat{y}$ i uporediti kvalitet tog preslikavanja za različite izbore funkcije reprezentacije $\phi(I) \mapsto \mathbf{x} \in \mathbb{R}^d$.
\end{itemize}

\subsection*{4.3 Posmatrane reprezentacije}

Reprezentacije se razlikuju po strukturnim osobinama koje su navedene u tabeli, i upravo se od tih razlika očekuje različito ponašanje pod različitim degradacijama:

\begin{center}
\begin{tabular}{|l|l|l|l|}
\hline
\textbf{Reprezentacija} & \textbf{Prostorni raspored} & \textbf{Domet} & \textbf{Učena?} \\
\hline
PCA & holistički, osetljiv na poravnanje & globalni & ne \\
\hline
HOG & kruta mreža, raspored očuvan & lokalni & ne \\
\hline
BoVW & \textbf{bez rasporeda}, raspored odbačen & lokalni & samo rečnik \\
\hline
CNN & hijerarhijski, invarijantan kroz \textit{pooling} & lokalni & u potpunosti \\
\hline
\end{tabular}
\end{center}

\textbf{Ključna metodološka odluka:} klasifikator se drži \emph{fiksnim} (linearni SVM) za sve četiri reprezentacije, kako bi razlike u rezultatima bile pripisive reprezentaciji, a ne klasifikatoru. Zbog toga se konvoluciona mreža evaluira \mvp{u dve varijante}: (a) kao ekstraktor obeležja, gde se izlaz pretposlednjeg sloja prosleđuje istom linearnom SVM klasifikatoru, čime postaje direktno uporediva sa ostale tri metode; i (b) end-to-end, sa sopstvenim softmax slojem. Ukupno se, dakle, poredi \textbf{pet konfiguracija}.

\subsection*{4.4 Očekivani rezultati}

\begin{itemize}
  \item \mvp{implementirane sve četiri reprezentacije i pet konfiguracija klasifikacije, nad jedinstvenim protokolom evaluacije};
  \item \mvp{podela na skupove za trening i validaciju izvršena po \textit{sekvencama} (eng. \textit{track}), a ne po pojedinačnim slikama} -- videti odeljak 6.1;
  \item \mvp{tabela glavnog poređenja}: tačnost, makro-F1 mera, vreme treniranja, vreme inferencije po slici, veličina modela i dimenzija vektora obeležja;
  \item \mvp{krive robusnosti} -- tačnost kao funkcija jačine degradacije, za četiri tipa degradacije;
  \item \mvp{analiza tačnosti u funkciji veličine znaka} -- test skup podeljen u opsege po visini graničnog okvira;
  \item \mvp{matrice konfuzije} za najbolju i najlošiju metodu, uz tabelu najčešće mešanih parova klasa;
  \item \mvp{tabela poređenja unapred zapisanih očekivanja sa dobijenim rezultatima}, sa obrazloženjem odstupanja;
  \item \str{ablaciona analiza pretprocesiranja} -- uticaj ekvalizacije histograma i izbora prostora boja;
  \item \str{pretraga hiperparametara} za svaku od metoda, umesto podrazumevanih vrednosti;
  \item \mvp{tehnički izveštaj sa analizom arhitekture rešenja i diskusijom rezultata}.
\end{itemize}

\newpage

\section*{5. Tehnologije koje će biti korišćene}

\begin{itemize}
  \item \textbf{Python} kao glavni programski jezik.
  \item \textbf{OpenCV}~\cite{opencv} za SIFT deskriptore, konverziju prostora boja, CLAHE ekvalizaciju i filtriranje.
  \item \textbf{scikit-image}~\cite{skimage} za HOG deskriptor -- pogodniji od \texttt{cv2.HOGDescriptor} za male kvadratne isečke, čije su podrazumevane vrednosti podešene za detekciju pešaka u prozoru $64\times128$.
  \item \textbf{scikit-learn}~\cite{sklearn} za PCA, \texttt{MiniBatchKMeans} (izgradnja vizuelnog rečnika), \texttt{LinearSVC} i metrike.
  \item \textbf{PyTorch}~\cite{pytorch} za konvolucionu mrežu.
  \item \textbf{NumPy, Pandas i Matplotlib} za obradu rezultata i generisanje grafika.
\end{itemize}

\textbf{Napomena o računskim resursima.} Razvojna mašina (AMD Ryzen 7 5800H, 8 jezgara / 16 niti, 13\,GB usable RAM) nema upotrebljiv GPU -- integrisana Vega grafika nije podržana od strane ROCm platforme, a CUDA nije primenljiva. Ceo projekat se stoga izvršava isključivo na procesoru. Ovo ograničenje diktira dve konkretne odluke: koristi se \texttt{LinearSVC} umesto SVM-a sa kernel funkcijom, čija je složenost $O(n^2)$--$O(n^3)$ i koja bi nad 39\,000 uzoraka trajala satima; i koristi se \texttt{MiniBatchKMeans} umesto klasičnog $k$-means algoritma pri izgradnji vizuelnog rečnika. Mala konvoluciona mreža (ispod milion parametara, ulaz $48\times48$) se na 16 niti trenira u prihvatljivom vremenu.

\section*{6. Validacija i podaci za testiranje}

\subsection*{6.1 Skup podataka i podela}

Koristi se \textbf{GTSRB} (\textit{German Traffic Sign Recognition Benchmark})~\cite{gtsrb} -- 43 klase, 39\,209 slika za trening i 12\,630 za testiranje, sa anotacijom koja za svaku sliku sadrži dimenzije, koordinate graničnog okvira znaka unutar isečka, i oznaku klase.

Skup ima osobinu koja se u praksi često previdi, tj. trening slike su grupisane u \textbf{sekvence od po 30 kadrova istog fizičkog znaka}, snimljenih dok mu se vozilo približava. Nasumična podela na trening i validaciju bi zato razdvojila skoro identične kadrove istog znaka na obe strane, što veštački podiže validacionu tačnost. \mvp{Podela se stoga vrši po identifikatoru sekvence, a ne po pojedinačnoj slici}, uz automatsku proveru da nijedna sekvenca ne pripada obema stranama podele.

Klase u GTSRB skupu su izrazito nebalansirane (odnos najzastupljenije i najređe klase je oko $10:1$), zbog čega se pored tačnosti obavezno izveštava i \mvp{makro-F1 mera}, koja svakoj klasi daje jednak značaj.

\subsection*{6.2 Kontrolisane degradacije}

Test skup se evaluira pod \textbf{tri} tipa kontrolisane degradacije, svaki na pet nivoa jačine:

\begin{enumerate}
  \item \mvp{\textbf{Aditivni Gausov šum}} -- $\sigma \in \{0, 5, 10, 20, 40\}$; simulira šum senzora pri slabom osvetljenju.
  \item \mvp{\textbf{Zamućenje usled kretanja}} -- dužina jezgra $\in \{0, 3, 5, 9, 15\}$ piksela, nasumičan ugao; simulira kretanje vozila.
  \item \mvp{\textbf{Promena osvetljenosti}} -- gama korekcija, $\gamma \in \{0.4, 0.7, 1.0, 1.5, 2.5\}$; simulira ekstremne uslove ekspozicije.
\end{enumerate}

Prvobitno je planirana i četvrta degradacija, perturbacija graničnog okvira, ali je analiza samog skupa podataka pokazala da ona nad GTSRB skupom \textbf{nije izvodljiva na pošten način}. GTSRB slike su već isečene oko znaka, sa medijanom margine od svega $16{,}7\%$ veličine okvira; sve izvan te margine odbačeno je pri formiranju skupa i ne postoji ni u jednoj datoteci. Jedini način da se takav pomeraj ipak proizvede jeste veštačko dopunjavanje ivica, čime bi se merio i postupak dopunjavanja, pored izabrane reprezentacije.

Ovakva vrsta degradacije se stoga izmešta na skupove podataka koji sadrže \textbf{kadrove punih scena}, a u ovom radu će se vršiti \str{evaluacija robusnosti opisanih modela nad GTSDB skupom pri nivoima pomeraja $\{0, 10, 20, 30, 40\}\%$ po razmeri i položaju}. GTSDB sadrži 1\,213 anotiranih znakova u 741 kadru dimenzija $1360\times800$, sa istom šemom od 43 klase, pa se pomeraj realizuje isecanjem iz stvarne scene, bez ijednog veštački generisanog piksela. Nulti nivo pomeraja reprodukuje kadriranje GTSRB skupa (okvir uvećan za oko $17\%$), kako bi razlika u odnosu na njega merila isključivo grešku okvira, a ne promenu kadriranja; razlika između čistog GTSRB rezultata i GTSDB rezultata pri nultom pomeraju zasebno se izveštava kao mera \emph{generalizacije između skupova podataka}. Zbog male i neravnomerne zastupljenosti klasa u GTSDB skupu (14 od 43 klase ima manje od 10 primeraka) na tom skupu se izveštava isključivo tačnost, ne i makro-F1.

Pomeraj okvira \textbf{ne ulazi u trening}, već se modeli treniraju isključivo nad GTSRB skupom, sa kadriranjem kakvo skup pruža. Time se meri \emph{inherentna} otpornost reprezentacije na grešku okvira, što je informativnije upravo zato što model takvu perturbaciju nikada nije video, a eksperimentalna kontrola ostaje neizmenjena. Ovo merenje je, kao što je objašnjeno u odeljku 4.1, direktna priprema za potencijalno buduće povezivanje sa modulom M1.% \footnote{Namerno se govori o sprezi M1$\to$M3, bez pominjanja modula M2, iako i on utiče na granični okvir koji stiže do klasifikatora. Razlog je u tome što merenje karakteriše osetljivost na \emph{grešku graničnog okvira} bez obzira na njen izvor, a M1 je primarni i neizbežni izvor te greške -- svaki sistem mora negde da locira znak, dok je M2 opciona optimizacija, budući da se detekcija može vršiti u svakom kadru. Uz to, M2 ne uvodi novu \emph{vrstu} greške, već doprinosi istoj: propagacijom okvira kroz kadrove akumulira se pomeraj (eng. \textit{drift}), koji se sabira sa greškom detektora. Veličina te akumulacije zavisi od izbora algoritma praćenja, učestanosti kadrova i intervala ponovne detekcije -- parametara koji u ovom projektu ne postoje i koji se ne mogu smisleno fiksirati. Kontekst se zato vezuje za M1, pri čemu dobijena vrednost predstavlja \emph{donju granicu} greške koju bi modul M3 video u kompletnom sistemu. \end{comment} }.

\subsection*{6.3 Metrike}

\begin{itemize}
  \item \textbf{Tačnost} i \textbf{makro-F1 mera} na test skupu;
  \item \textbf{F1 mera po klasi} i \textbf{matrica konfuzije};
  \item \textbf{tačnost po opsezima veličine znaka} -- opsezi $[0,32)$, $[32,48)$, $[48,72)$, $[72,\infty)$ piksela visine okvira; manji znak odgovara većoj udaljenosti u realnom scenariju;
  \item \textbf{krive robusnosti} -- tačnost u funkciji nivoa svake od tri degradacije, prikazana relativno u odnosu na tačnost datog metoda na čistom skupu, kako bi se poredila \emph{brzina opadanja} a ne polazna vrednost;
  \item \textbf{cena (\textit{computational price})} -- vreme treniranja, vreme inferencije po slici (medijana iz najmanje tri merenja, uz odbacivanje prvog poziva), veličina modela u megabajtima i dimenzija vektora obeležja.
\end{itemize}

\subsection*{6.4 Kriterijum uspešnosti}

Projekat se smatra uspešnim ako su sve četiri reprezentacije implementirane i evaluirane nad jedinstvenim protokolom, ako je podela po sekvencama sprovedena i proverena, i ako su dobijene krive robusnosti koje omogućavaju smislenu diskusiju o tome koja metoda otkazuje pod kojim uslovima. \textbf{Uspešnost se ne meri postignutom apsolutnom tačnošću}, jer je GTSRB zasićen skup na kome je visoka tačnost očekivana i sama po sebi ne predstavlja rezultat.

Pre pokretanja eksperimenata biće zapisana konkretna očekivanja o tome koja će reprezentacija biti najrobusnija pod kojom degradacijom, sa obrazloženjem izvedenim iz strukturnih osobina iz tabele u odeljku 4.3. Naknadno poređenje očekivanja sa rezultatima -- uključujući i objašnjenje onih očekivanja koja se nisu ostvarila -- predstavlja jezgro diskusije u izveštaju.

\section*{7. Literatura, linkovi i postojeća rešenja}

\bibliographystyle{unsrt}
\nocite{szeliski, milosavljevic, laganiere}
\bibliography{references}

\newpage

\subsection*{Dodatak A -- Razlozi odabira teme}

Tema je izabrana iz tri razloga.

\begin{enumerate}
	\item Prepoznavanje saobraćajnih znakova je zadatak koji prirodno objedinjuje veliki broj oblasti obrađenih na predmetu -- prostore boja i filtriranje kroz pretprocesiranje, detekciju karakterističnih tačaka i SIFT deskriptor kroz metod vreća vizuelnih reči, generativni pristup prepoznavanju kroz PCA (isti princip kao Eigenfaces metod), diskriminacioni pristup kroz HOG deskriptor, i konvolucione mreže kroz duboko učenje. Umesto da se svaka od ovih tema demonstrira izolovano, one se ovde stavljaju u direktno poređenje.

	\item Opseg projekta je namerno i eksplicitno ograničen na jedan modul šireg sistema. Autor smatra da je pošteno i korisnije realizovati jedan modul do kraja, sa urednom metodologijom i proverljivim rezultatima, nego površno dodirnuti ceo lanac obrade. Moduli koji su ostavljeni van opsega -- detekcija u punom kadru, praćenje kroz optički tok, vremenska agregacija odluka -- opisani su u odeljku 4.1 upravo zato da bi granica bila nedvosmislena.

	\item Projekat je zamišljen kao osnova za dalji rad. Klasifikator razvijen ovde postaje komponenta kompletnog sistema za prepoznavanje saobraćajne signalizacije iz video snimka, čiji bi naredni koraci bili: detekcija kandidatnih regiona nad GTSDB skupom, propagacija detekcija kroz kadrove pomoću optičkog toka~\cite{lucas-kanade} radi postizanja rada u realnom vremenu, i procena udaljenosti znaka na osnovu modela kamere i poznatih fizičkih dimenzija znaka. Merenje osetljivosti na pomeraj graničnog okvira, uvedeno u ovom projektu, direktno je motivisano tim nastavkom -- ono unapred kvantifikuje gubitak tačnosti pri prelasku sa anotiranih na detektovane isečke. Šire posmatrano, isti metodološki problem -- prepoznavanje objekata iz neprekidnog
toka kadrova snimljenih pokretnom kamerom iz prvog lica, uz stroga ograničenja
kašnjenja i računskih resursa -- predstavlja jedno od otvorenih pitanja egocentričnog
računarskog vida, oblasti koja poslednjih godina beleži izražen rast razvojem nosivih
uređaja sa kamerom; novija sveobuhvatna pregledna studija te oblasti izdvaja rad u
realnom vremenu i pri minimalnom kašnjenju kao zahtev koji postojeće metode još uvek
ne ispunjavaju~\cite{ego-survey}.
\end{enumerate}

\subsection*{Dodatak B -- Tabela akronima}

\begin{center}
\begin{tabular}{|l|p{10.5cm}|}
\hline
\textbf{Akronim} & \textbf{Značenje} \\
\hline
GTSRB & German Traffic Sign Recognition Benchmark \\
\hline
GTSDB & German Traffic Sign Detection Benchmark \\
\hline
PCA & Principal Component Analysis (analiza glavnih komponenti) \\
\hline
HOG & Histogram of Oriented Gradients (histogram orijentisanih gradijenata) \\
\hline
SIFT & Scale-Invariant Feature Transform \\
\hline
BoVW & Bag of Visual Words (vreća vizuelnih reči) \\
\hline
CNN & Convolutional Neural Network (konvoluciona neuronska mreža) \\
\hline
SVM & Support Vector Machine (metod potpornih vektora) \\
\hline
CLAHE & Contrast Limited Adaptive Histogram Equalization \\
\hline
ROI & Region of Interest (region od interesa) \\
\hline
RAM & Random Access Memory \\
\hline
GPU & Graphics Processing Unit \\
\hline
\end{tabular}
\end{center}

\end{document}
